% Generated by openai gpt-4.1 on 2026-11-09T13:06:33.868822+00:00, prompt geosmile-hull-v9 + card-mode addendum, temperature 0 (generate_card_narrative.py --stakeholder manager). Checks: passed. Do not edit by hand.
\renewcommand{\cardHeadline}{\kt{Crime history} dominates the model; \kt{Old Town} is the \kt{exception} in geographic sensitivity.}
\renewcommand{\cardSecA}{Panel 02 shows \kt{crime history} is the main driver across Hull, with census and built environment features next. \kt{East Ella} stands out: seasonality and unemployment lead there, not crime history. Only \pv{22} areas are not crime-led, and these rarely coincide with high geographic sensitivity. \kt{Old Town} is the clear exception for location sensitivity.}
\renewcommand{\cardSecB}{Inactive \%, road density, and nightlife count are among the \kt{features set aside}, each under \pv{1\%} and together \pv{8\%} of local importance. Related count and density features can split weight, so review them as groups. The most sensitive areas match \kt{expert regions} (see panel 01).}
\renewcommand{\cardFooter}{This shows what the model uses, not causes or resourcing priorities.}
\renewcommand{\cardModel}{gpt-4.1}
\renewcommand{\cardTileA}{\tile{145/168}{areas led by six-month crime average}}
\renewcommand{\cardTileB}{\tile{2.6x}{Old Town's sensitivity vs next area}}
\renewcommand{\cardTileC}{\tile{12/30}{features the model can set aside}}
\renewcommand{\cardTileD}{\tile{20/20}{top areas stable when neighbours added}}
